% cktdefs.tex

% THIS VERSION ASSUMES stdmath.tex IS INCLUDED

% MADE FROM THE MASTER SET OF DEFS FOR ALL OF EE101, 102, 103.
% math symbols and operators

\newcommand{\eqbydef}{\mathrel{\stackrel{\Delta}{=}}}
%\newcommand{\reals}{{\bf R}} % defined in stdmath.tex
%\newcommand{\ints}{{\bf Z}}
%\newcommand{\complex}{{\bf C}}
%\newcommand{\CN}{\complex^N}
\renewcommand{\Re}{\mbox{Re}}
\renewcommand{\Im}{\mbox{Im}}
\newcommand{\RMS}{\mbox{RMS}}

\renewcommand{\lg}{\mbox{lg}}

% some electrical units
\newcommand{\volt}{\mbox{\rm V}}
\newcommand{\microvolt}{\mbox{$\mu {\rm V}$}}
\newcommand{\millivolt}{\mbox{\rm mV}}
\newcommand{\kilovolt}{\mbox{\rm kV}}
\newcommand{\megavolt}{\mbox{\rm MV}}

\newcommand{\amp}{\mbox{\rm A}}
\newcommand{\picoamp}{\mbox {\rm pA}}
\newcommand{\nanoamp}{\mbox {\rm nA}}
\newcommand{\microamp}{\mu \mbox{\rm A}}
\newcommand{\milliamp}{\mbox{\rm mA}}
\newcommand{\kiloamp}{\mbox{\rm kA}}

\newcommand{\coul}{\mbox{coul}}

\newcommand{\ohm}{\mbox{$\Omega$}}
\newcommand{\megohm}{\mbox{{\rm M}$\Omega$}}
\newcommand{\kilohm}{\mbox{{\rm k}$\Omega$}}

\newcommand{\farad}{\mbox{\rm F}}
\newcommand{\picofarad}{\mbox{\rm pF}}
\newcommand{\nanofarad}{\mbox{\rm nF}}
\newcommand{\microfarad}{\mu \mbox{\rm F}}
\newcommand{\millifarad}{\mbox{\rm mF}}

\newcommand{\henry}{\mbox{\rm H}}
\newcommand{\microhenry}{\mu \mbox{\rm H}}
\newcommand{\millihenry}{\mbox{\rm mH}}

\newcommand{\watt}{\mbox{\rm W}}
\newcommand{\picowatt}{\mbox{\rm pW}}
\newcommand{\microwatt}{\mu \mbox{\rm W}}
\newcommand{\milliwatt}{\mbox{\rm mW}}
\newcommand{\kilowatt}{\mbox{\rm kW}}
\newcommand{\megawatt}{\mbox{\rm MW}}
\newcommand{\gigawatt}{\mbox{\rm GW}}

\newcommand{\gram}{\mbox{\rm g}}
\newcommand{\kilogram}{\mbox{\rm kg}}
\newcommand{\newton}{\mbox{\rm N}}
\newcommand{\kilonewton}{\mbox{\rm kN}}
\newcommand{\joule}{\mbox{\rm J}}
\newcommand{\kilojoule}{\mbox{\rm kJ}}

\newcommand{\degK}{\mbox{${}^{\circ}$\rm K}}
\newcommand{\degC}{\mbox{${}^{\circ}$\rm C}}
\newcommand{\degF}{\mbox{${}^{\circ}$\rm F}}

\newcommand{\foot}{\mbox{\rm ft}}

\newcommand{\mins}{\mbox{min}} %to prevent confusion with \min, the minimum operator
\newcommand{\maxs}{\mbox{max}}

\renewcommand{\sec}{\mbox{sec}}
\newcommand{\millisec}{\mbox{msec}}
\newcommand{\microsec}{\mu\mbox{sec}}
\newcommand{\nanosec}{\mbox{nsec}}

\newcommand{\rad}{\mbox{rad}}
\newcommand{\hertz}{\mbox{Hz}}
\newcommand{\kilohertz}{\mbox{kHz}}
\newcommand{\megahertz}{\mbox{MHz}}
\newcommand{\dB}{\mbox{dB}}


\newcommand{\rms}{\mbox{RMS}}
\newcommand{\avg}{\mbox{AVG}}

\newcommand{\eqmaybe}{\mathrel{\stackrel{?}{=}}}

% others moved to stdmath.h
\newcommand{\BEAL}{\begin{eqnarray}}
\newcommand{\EEAL}{\end{eqnarray}}

\newcommand{\XXX}[1]{\marginpar{\raggedright \footnotesize #1}}

\newtheorem{remark}{Remark}


% Usage: \rawFigure{figure.eps}
%\def\figureskip{\parskip}
\def\figureskip{0pt} % deferred evaluation
\def\rawFigure#1{%
    \vskip\figureskip\nobreak
    \hbox to \hsize{\epsfbox{#1}\hfil}
    \vskip\figureskip
}

\setlength{\parskip}{2ex plus 0.5ex minus 0.5ex }

% Call this after \maketitle, and possibly also after any outline and after
% \makenavtoc (stdcommon.tex)
%begin{latexonly}
\newcommand{\setlecfont}{
  \sffamily\LARGE
  \raggedright
}
%end{latexonly}

\begin{htmlonly}
\newcommand{\setlecfont}{
% Creates large italics for some reason in HTML case.
}
\end{htmlonly}

%begin{latexonly}
\newcommand{\oursection}[1]{
\clearpage
\begin{center}
{\Huge\bfseries #1}
\end{center}
\vspace*{0.15cm}
\hrule height.3mm
\bigskip
}
\newcommand{\oursectionncp}[1]{
\begin{center}
{\Huge\bfseries #1}
\end{center}
\vspace*{0.15cm}
\hrule height.3mm
\bigskip
}
%end{latexonly}

\begin{htmlonly}
\newcommand{\oursection}[1]{
%7/28/01/jos: Changed ``renew'' to ``new'' above (to clear images.log error)
%4/9/01/jos: ``huge'' is now coming out as H1+4 --- what changed?
%4/9/01/jos/later: I think I needed ``renew'' above.  LARGE is now H1+2
\section*[#1]{{\LARGE\bfseries #1}} % LARGE formerly subsection heading size
%to 4/9/01: \section*[#1]{{\huge\bfseries #1}} % NOTE DIFFERENT FONT SIZE (+3 not +4)
%\section*[#1]{\begin{center}{\Huge\bfseries #1}\end{center}}
%\begin{rawhtml}
%<HR width="50\%">
%\end{rawhtml}
}
\end{htmlonly}

%begin{latexonly}
\newcommand{\oursubsection}[1]{
\begin{center}
{\LARGE\bfseries #1}
\end{center}
%\vspace*{0.15cm}
\bigskip
}
\newcommand{\oursubsectionsamepage}[1]{\oursubsection{#1}}
\newcommand{\oursubsubsection}[1]{\subsubsection*{{\LARGE\bfseries #1}}
}
%end{latexonly}

\begin{htmlonly}
\newcommand{\oursubsection}[1]{\subsection*[#1]{{\LARGE\bfseries #1}}}
\newcommand{\oursubsectionsamepage}[1]{{\LARGE\bfseries #1}}
\newcommand{\oursubsubsection}[1]{{\LARGE\bfseries #1}}
%Makes new pages in l2h - could fix with MAX_SPLIT_DEPTH:
% \subsubsection*{{\LARGE\bfseries #1}}
\end{htmlonly}

\newcommand{\figureCenterCapSize}[4]{
	\begin{center}
%	\hline
	\epsfxsize=#4
	\leavevmode
	\epsfbox{#2}
%	\LARGE #3
%	\hline
	\end{center}
}

\newcommand{\figureCenterCap}[3]{
	\figureCenterCapSize{#1}{#2}{#3}{7in}
}

\newcommand{\figureLongCap}[3]{
	\figureCenterCapSize{#1}{#2}{#3}{7in}
}

\newcommand{\mysection}[1]{\oursection{#1}\message{WARNING: l2h will NOT GENERATE A NEW PAGE}}
\newcommand{\mysubsection}[1]{\oursubsection{#1}\message{WARNING: l2h will NOT GENERATE A NEW PAGE}}
\newcommand{\mysubsubsection}{\subsubsection}

% Also in stdmath.tex
\providecommand{\zbox}[1]{\ensuremath{\fbox{$\displaystyle #1$}}}

\newcommand{\xlabelsize}{\footnotesize}
\newcommand{\mylabelsize}{\footnotesize}
\newcommand{\leclabelsize}{\large}
\newcommand{\psfragsize}{\large}

\newcommand{\fs}{f_s}
